\begin{titlepage}
\keepXColumns % Verhindert, dass ltablex die Spaltenbreiten falsch berechnet

\vspace*{1cm}

\noindent Report nach IPMA ICB 4.0\\[.5cm]
\textbf{{\Large IPMA\textsuperscript{\textregistered}{} LEVEL B - Certified Senior Project Manager}}\\[0.5cm]

\setlength{\extrarowheight}{2pt}
\renewcommand{\arraystretch}{1.5}

% Da ltablex tabularx zu einer longtable macht, nutzen wir \noindent 
% und setzen das optionale Argument [l] für linksbündig.
\noindent
\begin{tabularx}{\textwidth}{|l|X|}
\hline 
Name, Vorname                                            & Seitz, Sebastian \\ \hline 
Titel des Dokuments:                                     & Report\_Seitz\_Sebastian\_19871112\_V1.pdf \\ \hline 
Eingereicht zur Zertifizierungsrunde IZR/ OZRD           & \curse \\ \hline 
Datum:                                                   & \today {} {} \timestamp {} \\ \hline 
Version:                                                 & 1.0 \\ \hline 
\end{tabularx}

\vspace*{1cm}

\noindent
\begin{tabularx}{\textwidth}{|l|l|X|l|}
\hline
\rowcolor[HTML]{9B9B9B} 
\bfseries Version & \bfseries Datum & \bfseries Änderung & \bfseries Seite/n \\ \hline
1.0 & \ersterImMonat{} & Erstellung des Reports & alle \\ \hline
\end{tabularx}

\vspace*{0.5cm}
\noindent München, den \today

\end{titlepage}